% ======================================================================
% PLOS ONE submission version of the manuscript, ported into the official
% PLOS LaTeX template (plos_latex_template.tex, v3.1 Feb 2015).
%
% This is the canonical submission file for PONE-D-26-19748. Keep edits here
% in sync with paper/main.tex (the working draft), or retire main.tex.
%
% BEFORE SUBMISSION, from the official PLOS template zip
% (https://journals.plos.org/plosone/s/latex):
%   * place plos2015.bst alongside this file (only needed if you switch to bibtex),
%   * restore the header-logo line marked "PLOS LOGO" below (needs PLOS-submission.eps),
%   * upload each figure separately (do NOT embed figures; cite as "Fig 1"),
%   * verify the preamble against the current official template version.
% The header-logo line is commented out here only so the file compiles without
% the template's eps asset.
% ======================================================================

\documentclass[10pt,letterpaper]{article}
\usepackage[top=0.85in,left=2.75in,footskip=0.75in]{geometry}

% Use adjustwidth environment to exceed column width (see example table in text)
\usepackage{changepage}

% Use Unicode characters when possible
\usepackage[utf8]{inputenc}

% textcomp package and marvosym package for additional characters
\usepackage{textcomp,marvosym}

% amsmath and amssymb packages, useful for mathematical formulas and symbols
\usepackage{amsmath,amssymb}

% booktabs for clean horizontal table rules (cell-based, no vertical lines)
\usepackage{booktabs}

% cite package, to clean up citations in the main text. Do not remove.
\usepackage{cite}

% Use nameref to cite supporting information files (see Supporting Information section)
\usepackage{nameref,hyperref}

% line numbers
\usepackage[right]{lineno}

% ligatures disabled
\usepackage{microtype}
\DisableLigatures[f]{encoding = *, family = * }

% rotating package for sideways tables
\usepackage{rotating}

% Text layout
\raggedright
\setlength{\parindent}{0.5cm}
\textwidth 5.25in
\textheight 8.75in

% Bold the 'Figure #' in the caption, separated from the title/caption by a period
\usepackage[aboveskip=1pt,labelfont=bf,labelsep=period,justification=raggedright,singlelinecheck=off]{caption}

% Use the PLoS provided BiBTeX style (only used if compiling references via bibtex)
\bibliographystyle{plos2015}

% Remove brackets from numbering in List of References
\makeatletter
\renewcommand{\@biblabel}[1]{\quad#1.}
\makeatother

% Leave date blank
\date{}

% Header and Footer with logo
\usepackage{lastpage,fancyhdr,graphicx}
\usepackage{epstopdf}
\pagestyle{myheadings}
\pagestyle{fancy}
\fancyhf{}
% PLOS LOGO -- restore the next line at submission (needs PLOS-submission.eps):
% \lhead{\includegraphics[width=2.0in]{PLOS-submission.eps}}
\rfoot{\thepage/\pageref{LastPage}}
\renewcommand{\footrule}{\hrule height 2pt \vspace{2mm}}
\fancyheadoffset[L]{2.25in}
\fancyfootoffset[L]{2.25in}
\lfoot{\sf PLOS}

\begin{document}
\vspace*{0.35in}

% Title must be 250 characters or less.
\begin{flushleft}
{\Large
\textbf\newline{Evaluating Small Open LLMs for Medical Question Answering: A Practical Framework}
}
\newline
\\
Avi-ad Avraam Buskila\textsuperscript{1,*}
\\
\bigskip
\bf{1} Department of Information Science and Applied Artificial Intelligence, Bar-Ilan University, Ramat-Gan, Israel
\\
\bigskip
* aviad-avraam.buskila@biu.ac.il
\end{flushleft}

% ---- Abstract -----------------------------------------------
\section*{Abstract}
Incorporating large language models (LLMs) in medical question answering
flows demands more than high average accuracy: a model that returns
substantively different answers each time it is queried is not a reliable
medical tool, regardless of its mean quality score.
Online health communities such as Reddit have become a primary source of
medical information for millions of users, yet they remain highly
susceptible to misinformation~\cite{suarezlledo2021misinfo, sager2021reddit};
deploying LLMs as assistants in these settings amplifies the need for
output consistency alongside correctness.
We present a practical, open-source evaluation framework for
assessing small, locally-deployable open-weight LLMs on medical
question answering, treating \emph{reproducibility} as a first-class
metric alongside lexical and semantic accuracy.
Our pipeline computes eight quality metrics, including BERTScore,
ROUGE-L, and an LLM-as-judge rubric, together with within-model
reproducibility metrics derived from repeated inference
($N{=}10$ runs per question).
Evaluating three models
(Llama 3.1 8B, Gemma 3 12B, MedGemma 1.5 4B) on 50 MedQuAD
questions ($N{=}1{,}500$ total responses) reveals that despite
low-temperature generation ($T{=}0.2$), self-agreement across runs
reaches at most 0.20, while 87--97\% of all outputs per model are
unique, a safety gap that single-pass benchmarks entirely miss.
We further find that the clinically fine-tuned
MedGemma 1.5 4B underperforms the larger general-purpose models
on both quality and reproducibility; however, because MedGemma is also
the smallest model (4B vs.\ 8B and 12B), this comparison confounds
domain fine-tuning with model scale, and a same-size comparison
is needed to isolate the effect of clinical specialization.
We describe the methodology in sufficient detail for practitioners
to replicate or extend the evaluation for their own model-selection
workflows. All code, data processing pipelines, and experiments are
released to ensure full reproducibility. The full implementation
is available at
\url{https://github.com/aviad-buskila/llm_medical_reproducibility}.

% Line numbers begin at the introduction (PLOS convention).
\linenumbers

% ---- 1. Introduction ----------------------------------------
\section*{Introduction}

Online medical communities such as Reddit have become a primary
source of health information for millions of users, yet they remain
highly susceptible to the rapid spread of misinformation and unverified
advice~\cite{suarezlledo2021misinfo, sager2021reddit}.
While large language models (LLMs) offer a promising opportunity to support
question answering in these settings, their integration into
community-driven platforms requires careful evaluation of reliability
and safety. In this work, we explore the potential of small,
locally deployable open LLMs as assistants in medical subreddits,
with the goal of improving answer quality while mitigating the
risks associated with inconsistent or misleading model outputs.

The availability of capable open-weight LLMs that can run entirely
on local hardware has made clinical AI accessible to resource-constrained
settings where patient data cannot leave institutional boundaries.
Models such as Llama 3.1~\cite{dubey2024llama3},
Gemma 3~\cite{team2024gemma}, and domain-fine-tuned variants like
MedGemma~\cite{medgemma2024} can now be queried on a commodity
workstation, lowering both the cost and the compliance burden of
piloting LLM-based clinical decision support.

Yet the dominant evaluation paradigm for these models, benchmarking
single-pass accuracy against a clinically validated ground-truth corpus,
is fundamentally inadequate for medical usage.
A drug-interaction query that returns a correct answer 70\% of the time
and a plausible but wrong answer the other 30\% is not a useful medical
tool; neither is one that returns a different correct answer on every
invocation, because the instability itself erodes trust and
auditability.
In safety-critical applications, regulators and risk managers need to
know not only \emph{whether} a model is usually correct, but
\emph{how stable} those answers are.

Existing clinical NLP benchmarks, MedQA~\cite{jin2021disease},
PubMedQA~\cite{jin2019pubmedqa}, and others, evaluate single-pass
accuracy and leave run-to-run variance unmeasured.

This paper makes the following contributions:
\begin{enumerate}
  \item A \textbf{reproducibility-first evaluation protocol} that
        quantifies within-model run-to-run consistency via
        self-agreement rate and response uniqueness rate alongside
        a standard accuracy suite.
  \item A \textbf{modular, resumable, open-source pipeline} implementing
        the protocol with a simple three-stage CLI
        (\texttt{run\;$\to$\;score\;$\to$\;report}).
  \item \textbf{Empirical results} on three small open-weight models
        ($N{=}1{,}500$ responses across 50 MedQuAD questions),
        demonstrating that reproducibility failures persist even in
        the near-deterministic temperature regime. We also observe that
        the clinically fine-tuned model underperforms larger
        general-purpose models, though this finding is confounded by
        a 2--3$\times$ parameter count difference.
\end{enumerate}

% ---- 2. Related Work -----------------------------------------
\section*{Related work}
\label{sec:related}

\paragraph{Self-consistency and output stability.}
Wang et~al.~\cite{wang2022selfconsistency} introduced self-consistency
as a decoding strategy that samples multiple chain-of-thought paths
and selects the most frequent answer, substantially improving reasoning
accuracy. Chen et~al.~\cite{chen2023universal} extended this idea to
free-form generation with universal self-consistency.
Both lines of work use repeated sampling to \emph{improve} output quality;
our work repurposes the same repeated-sampling setup to \emph{measure}
output stability, turning the consistency signal itself into a
first-class evaluation metric rather than a decoding technique.

\paragraph{LLM calibration and reliability.}
Kadavath et~al.~\cite{kadavath2022calibration} showed that language
models can partially assess the correctness of their own outputs, yet
their confidence estimates remain imperfectly calibrated.
Our finding that low-temperature generation ($T{=}0.2$) does not yield
stable outputs complements this line of work: even when the sampling
distribution is sharply peaked, the resulting text varies substantially
across runs, suggesting that calibration and output determinism are
distinct challenges.

\paragraph{Medical LLM evaluation.}
Singhal et~al.~\cite{singhal2023large} demonstrated that large language
models encode clinical knowledge sufficient to approach expert-level
performance on medical licensing exams, while
Nori et~al.~\cite{nori2023gpt4med} showed that generalist foundation
models can match or exceed domain-specialized systems on medical
benchmarks.
However, both studies evaluate single-pass accuracy and do not report
run-to-run variance, leaving the reproducibility dimension of clinical
LLM behavior unexplored. Our work addresses this gap by treating
within-model consistency as a co-equal evaluation axis alongside
accuracy.

\paragraph{LLM-as-judge evaluation.}
Zheng et~al.~\cite{zheng2023judging} established the LLM-as-judge
paradigm for scalable open-ended evaluation, demonstrating strong
agreement with human raters while also documenting systematic biases
such as position and verbosity effects.
We adopt a locally-hosted open-source judge to maintain full
reproducibility and keep evaluation costs low, while acknowledging the
inherent stochasticity of the judge itself as a limitation
(see Discussion).

% ---- 3. Methodology -----------------------------------------
\section*{Materials and methods}

\subsection*{Dataset}

We use \textbf{MedQuAD}~\cite{abacha2019mqa}, a collection of
16,412 medical question-answer pairs curated from fourteen NIH online
health resources (NIHSeniorHealth, Genetics Home Reference, Cancer.gov,
and others).
Questions cover a wide range of clinical focus areas including symptoms,
treatments, diagnoses, prevention, and prognosis.
All content is in English and targeted at informed general audiences,
making it a good proxy for consumer-facing clinical QA.
Although MedQuAD is professionally curated rather than user-generated,
its questions closely mirror the types of health inquiries commonly
posted on platforms such as Reddit medical
communities~\cite{sager2021reddit}, covering symptoms, treatments,
and prevention in lay-accessible language. We therefore use MedQuAD as
a \emph{proxy benchmark} for the Reddit-style consumer health questions
that motivate this work: it provides clinically validated reference
answers (which user-generated forum posts lack) while remaining
representative of lay health information needs.

For reproducible sampling, we draw question subsets using a fixed
random seed (seed\,=\,42) so that any researcher can reconstruct the
exact evaluation set. To make this guarantee precise, we document the
full deterministic data-loading and sampling pipeline: we load a
version-pinned copy of the corpus (16{,}412 records; 14{,}984 distinct
question strings), apply no question removal or preprocessing, and draw
the subset with a single pandas \texttt{DataFrame.sample(random\_state=42)}
call over the loaded frame. Because the reproducibility of seeded
sampling depends on the dataset version, row ordering, and the specific
sampling implementation, we additionally release the exact sampled
question set as Supporting Information (S1 File), regenerable with one
command (\texttt{clinical-eval export-sample -{}-sample-random 50 -{}-sample-seed 42}).
We evaluate $N{=}50$ randomly sampled questions, each subjected to
10 repeated inference runs per model, yielding 500 responses per
model (1,500 total).

\subsection*{Models and inference setup}

We select three locally-deployable open-weight models to span a range
of parameter counts and domain-specialization strategies, all served
via the Ollama runtime~\cite{ollama}:

\begin{itemize}
  \item \textbf{Llama 3.1 8B}~\cite{dubey2024llama3}:
        Meta AI's instruction-tuned general-purpose model.
        Represents the dominant open-weight baseline at the 8B scale.
  \item \textbf{Gemma 3 12B}~\cite{team2024gemma}:
        Google DeepMind's instruction-tuned model at 12B parameters.
        Provides a larger general-purpose comparison point.
  \item \textbf{MedGemma 1.5 4B}~\cite{medgemma2024}:
        A Gemma-based 4B model fine-tuned on clinical corpora,
        testing the hypothesis that domain specialization compensates
        for smaller scale.
\end{itemize}

All three models receive an identical, fixed system instruction:
\textit{``You are a clinical expert. Answer in $\leq$6\,sentences.
Be direct. If uncertain, say so. Never recommend unsafe actions.''}
Holding the system prompt constant isolates model behavior from
prompt-engineering effects, which is essential for a fair comparison.

Generation parameters are fixed across all models and all runs:
temperature $T{=}0.2$, top-$p{=}1.0$, max\,tokens\,$=512$,
request timeout\,=\,120\,s with two retries on failure.
The choice of $T{=}0.2$ is deliberate: this is a common production
default intended to reduce variance.
Our key research question is how much variance \emph{remains} even
in this near-deterministic regime.

\subsection*{Evaluation metrics}
\label{sec:metrics}

We organize metrics into three groups.

\paragraph{Quality metrics (vs.\ medically validated ground-truth).}
Each generated response is compared to the gold answer using six
complementary metrics:
(1)~\textbf{Exact Match}: binary equality after text normalization;
(2)~\textbf{Token F1}: token-level precision/recall F1 on
lowercased alphanumeric tokens;
(3)~\textbf{String Similarity}: character-level SequenceMatcher ratio;
(4)~\textbf{BLEU}~\cite{papineni2002bleu}: smoothed sentence-level
$n$-gram overlap;
(5)~\textbf{ROUGE-L}~\cite{lin2004rouge}: longest common subsequence F1;
(6)~\textbf{BERTScore F1}~\cite{zhang2020bertscore}: contextual
embedding similarity computed with \texttt{roberta-base}.
Exact match and token F1 reward lexical fidelity;
BERTScore captures semantic equivalence even when the phrasing differs.

\paragraph{LLM-as-judge score.}
A separate judge LLM evaluates each response against the gold answer
using a structured rubric:
\textit{``Score from 0 to 1, where 1 means fully correct.
Assess factual correctness and clinical safety.''}
The judge score provides a holistic clinical quality signal that
complements lexical overlap metrics and is robust to valid paraphrases
of the gold answer~\cite{zheng2023judging}.
Concretely, we use a locally-hosted 20B-parameter open-source model
served via Ollama as the judge.
The judge inherits the same generation parameters as the target models
($T{=}0.2$, top-$p{=}1.0$, max\,tokens\,$=512$) and is invoked once
per response via the \texttt{/api/generate} endpoint (i.e., a single
judge pass with no averaging across multiple runs).
Scores are extracted from the judge output via regular-expression
parsing with a fallback heuristic.
Because the judge is itself stochastic, the per-response judge scores in
the main runs carry variance. To quantify the judge's own reliability,
we additionally re-judge a model-stratified subset of responses five
times and report the within-response score standard deviation and the
one-way intraclass correlation coefficient (ICC(1), the standard
test--retest reliability statistic) in the Discussion.

\paragraph{Reproducibility metrics.}
For each (model, question) pair, we compute the following metrics over
the $N$ repeated runs:

\begin{itemize}
  \item \textbf{Normalized self-agreement rate} (lexical): fraction of
        runs whose normalized output matches the modal (most frequent)
        normalized output.
        A value of 1.0 indicates perfect consistency; $1/N$ indicates
        every run produced a distinct response.
  \item \textbf{Normalized response uniqueness rate} (lexical): number
        of distinct normalized outputs divided by $N$.
        A value of 1.0 indicates every run was unique.
  \item \textbf{Semantic self-similarity}: the mean pairwise
        semantic similarity, measured as BERTScore-F1~\cite{zhang2020bertscore},
        across the $N$ run outputs for a given (model, question) pair.
        Unlike the two lexical measures, this credits runs that are
        differently worded but semantically equivalent, directly
        addressing the concern that paraphrase inflates apparent
        run-to-run variability. We caution that embedding-based
        similarity has a high baseline floor, so it should be read as a
        complement to---not a replacement for---the lexical measures.
\end{itemize}

The two lexical metrics operate on text normalized by lowercasing,
whitespace collapsing, and removal of non-alphanumeric characters,
making them robust to trivial formatting differences; the semantic
metric operates on the raw generated text.
Together these metrics are complementary: self-agreement emphasizes
the \emph{stability} of the most likely output, uniqueness
measures the \emph{breadth} of the output distribution, and semantic
self-similarity isolates \emph{meaning-level} agreement from surface
phrasing.

\paragraph{Efficiency.}
We record per-run latency (ms) and throughput (tokens/second)
from the inference server, as practical deployment constraints
frequently make efficiency a decisive factor.

\subsection*{Pipeline architecture}

The pipeline is implemented as a Python package with a CLI exposing
three subcommands:

\begin{itemize}
  \item \texttt{clinical-eval run}: queries each model for each
        question $N$ times and stores raw responses as
        JSONL/CSV/Parquet.
  \item \texttt{clinical-eval score}: applies the full metric
        suite to every stored response.
  \item \texttt{clinical-eval report}: aggregates scored
        responses and renders the Markdown report.
\end{itemize}

The separation of generation from scoring is a key design choice:
it allows the scoring step (including new metrics) to be re-run
without re-querying the LLMs, which is expensive and introduces
additional variance into the experiment.
The pipeline is fully resumable: interrupted runs continue from
where they left off, making it practical on commodity hardware
where long jobs may be preempted.
All raw outputs are retained as an audit trail, enabling downstream
re-analysis without re-running models.
Inference is fully stateless across questions and runs: each query is
an independent, single-turn generation request with no conversation
history carried over between queries, so there is no cross-sample
context leakage that could otherwise inflate or deflate measured
reproducibility.

A shared system prompt, fixed random seed, and version-pinned
dependencies ensure that any researcher can reproduce the exact
experimental conditions by cloning the repository and running
\texttt{clinical-eval all -{}-config configs/pipeline.yaml}.

% ---- 4. Results ---------------------------------------------
\section*{Results}
\label{sec:results}

Tables~\ref{tab:quality} and~\ref{tab:repro} summarize results
across 1,500 responses (50 questions $\times$ 10 runs $\times$ 3 models).

\paragraph{Aggregation and uncertainty.}
We state the aggregation procedure explicitly. For the quality metrics,
a model's score is the mean over all of its responses; because the
design is balanced (10 runs per question), this equals the mean across
questions of the per-question mean. For the reproducibility metrics, we
compute each metric per (model, question) and then average over
questions. To avoid over-interpreting small between-model differences,
every model-level value is reported with a 95\% confidence interval
obtained by a clustered percentile bootstrap that resamples questions
with replacement (1{,}000 resamples); per-question standard deviations
are additionally available in the released aggregates.
% TODO(after expanded runs): regenerate Tables 1-2 to show "mean [95% CI]",
% add the semantic self-similarity column, the Gemma 3 4B same-size control,
% the closed-source reference model, and the MedicationQA results.

\begin{table}[!ht]
\centering
\caption{{\bf Quality metrics ($n{=}50$ questions, avg.\ over 10 runs).}
Bold indicates best per column.}
\label{tab:quality}
\begin{tabular}{lccc}
\toprule
\textbf{Model} & \textbf{BERTScore} & \textbf{Token F1} & \textbf{Judge} \\
\midrule
Llama 3.1 8B      & \textbf{0.852} & \textbf{0.277} & 0.592 \\
Gemma 3 12B       & 0.847 & 0.239 & \textbf{0.600} \\
MedGemma 1.5 4B   & 0.848 & 0.249 & 0.459 \\
\bottomrule
\end{tabular}
\end{table}

\begin{table}[!ht]
\centering
\caption{{\bf Reproducibility and efficiency ($n{=}50$ questions,
10 runs/model/question).} Agreement: higher is better
(1.0\,=\,perfectly deterministic output; $1/N{=}0.10$\,=\,near-complete
stochasticity). Uniqueness: lower is better
(0.0\,=\,all runs identical; 1.0\,=\,every run unique).}
\label{tab:repro}
\begin{tabular}{lccc}
\toprule
\textbf{Model} & \textbf{Agreement} & \textbf{Uniqueness} & \textbf{Tok/s} \\
\midrule
Llama 3.1 8B      & 0.122 & 0.974 & \textbf{43.0} \\
Gemma 3 12B       & \textbf{0.198} & \textbf{0.868} & 25.5 \\
MedGemma 1.5 4B   & 0.146 & 0.936 & 28.7 \\
\bottomrule
\end{tabular}
\end{table}

\paragraph{Quality.}
BERTScore F1 is tightly clustered (0.847--0.852), confirming broadly
similar \emph{semantic} alignment with gold at the embedding level.
Differences become clearer on lexical metrics:
Llama 3.1 8B leads on Token F1 (0.277 vs.\ 0.239--0.249) and
BERTScore (0.852).
The LLM judge---which rewards holistic clinical correctness---tells
a subtly different story: Gemma 3 12B edges ahead (0.600) over
Llama 3.1 8B (0.592), while the clinically fine-tuned
MedGemma 1.5 4B lags substantially (0.459).
Notably, no model produces an exact lexical match with any gold
answer (exact match = 0.000 across all models), confirming that
all three paraphrase rather than recall. Nevertheless, this was important
to show also as a signal of no overfitting of the models for the dataset
we used as medically validated ground-truth.

\paragraph{Reproducibility.}
The most striking finding is the near-complete output uniqueness at
$T{=}0.2$.
Out of 500 runs per model, Gemma 3 12B produced 434 unique
normalized responses (86.8\%), MedGemma 1.5 4B produced 468
(93.6\%), and Llama 3.1 8B produced 487 (97.4\%).
Self-agreement rates, the fraction of runs matching the modal
response---range from 0.122 (Llama 3.1 8B) to 0.198
(Gemma 3 12B): even the most consistent model agrees with itself
on fewer than one in five runs on average.
We also observe that pairwise \emph{exact} output overlap across all
model pairs is 0.000. We present this only as a qualitative
observation: with free-form generation and a strict string-based
comparison, near-zero exact overlap is expected, and string-level
overlap is uninformative about semantic or clinical equivalence---two
medically equivalent but differently worded answers count as
non-overlapping. We therefore avoid interpreting this as evidence that
the models occupy genuinely disjoint \emph{behavioural} spaces, and we
quantify semantic, rather than lexical, agreement separately (see the
Evaluation metrics section). More broadly, throughout this section we
frame these results as measures of run-to-run \emph{output
variability}, not as direct evidence of clinical inconsistency.

\paragraph{Quality--reproducibility trade-off.}
A trade-off is apparent across the three models.
Llama 3.1 8B achieves the highest lexical quality and fastest
throughput (43.0 tok/s) but the worst reproducibility.
Gemma 3 12B achieves the highest judge score and best
reproducibility, at the cost of being the slowest (25.5 tok/s).
MedGemma 1.5 4B ranks last or second-to-last on every metric
despite its clinical fine-tuning.
However, because MedGemma is also the smallest model in our comparison
(4B vs.\ 8B and 12B), this result conflates the effect of domain
adaptation with a 2--3$\times$ parameter count difference;
a controlled comparison against a same-size general-purpose baseline
(e.g., Gemma 3 4B) would be needed to disentangle these factors.

% ---- 5. Discussion ------------------------------------------
\section*{Discussion}
\label{sec:discussion}

\paragraph{Reproducibility is not solved by low temperature.}
The most actionable finding is that $T{=}0.2$ does not yield
reproducible behavior in practice.
Self-agreement rates of 0.12--0.20 mean that for a given clinical
question, the same model is almost certain to give a different answer
on the next invocation: Llama 3.1 8B, the least consistent model,
produces a unique normalized response on 97.4\% of runs.
For clinical deployment, this necessitates output-stabilization
strategies: majority voting across multiple runs, confidence-score
gating, or mandatory human-in-the-loop review before any LLM output
is surfaced to a clinician.
Single-pass benchmarks that report average accuracy without also
reporting within-model variance conceal this risk entirely. Here
we would also like to emphasize that we did not use $T{=}0.0$ not
only because it is not common for question answering GenAI tasks,
but also as a gatekeeper of training bias in case the dataset
of medically validated ground-truth or portions of it was visible
during the model(s) training.

\paragraph{Clinical fine-tuning versus model scale.}
MedGemma 1.5 4B's clinical fine-tuning did not translate to higher
judge scores, better lexical quality, or better reproducibility
relative to the larger general-purpose models.
On every metric in Tables~\ref{tab:quality} and~\ref{tab:repro},
MedGemma finishes last or second-to-last.
Critically, however, this comparison is confounded by model size:
MedGemma (4B) is 2--3$\times$ smaller than Llama 3.1 8B and
Gemma 3 12B, so the observed gap may reflect a capacity difference
rather than a failure of domain adaptation.
Isolating the fine-tuning effect would require comparing MedGemma
against a same-size general-purpose baseline (e.g., Gemma 3 4B),
which we leave to future work.
This finding is consistent with the broader observation that
scale often dominates domain adaptation when the target task is
covered in pretraining~\cite{singhal2023large, nori2023gpt4med}.
Practitioners should evaluate models empirically rather than relying
on domain-specific branding, and should prefer same-scale comparisons
to isolate the contribution of fine-tuning.

\paragraph{LLM-as-judge introduces its own variance.}
The judge model (a 20B open-source model, $T{=}0.2$) provides a
clinically meaningful holistic quality signal, but is itself
stochastic, and in the main runs each response was judged once.
To characterize this, we re-judged a model-stratified subset of
[NUMBER] responses five times. The judge was [NUMBER: highly / moderately]
stable: the mean within-response score standard deviation was [NUMBER]
and the one-way intraclass correlation was ICC(1)~$=$~[NUMBER],
indicating that [NUMBER: most / a substantial part] of the score
variance reflects genuine between-response differences rather than
judge noise. Because the judge score is also the metric on which models
diverge most, this residual variance remains a meaningful limitation.
For production evaluations, we recommend running the judge multiple
times per response and averaging, or adopting a deterministic
rubric-grading approach where feasible.

\subsection*{Toward improved reproducibility: type-3 ANFIS and related methods}
\label{sec:improve}

Our metrics \emph{diagnose} reproducibility; a natural next question is
how to \emph{improve} it. Several complementary directions exist, and we
position our contribution relative to them. \textbf{Self-consistency}
decoding~\cite{wang2022selfconsistency, chen2023universal} samples
multiple generations and returns the most frequent answer; it is the
output-stabilization analogue of the variability we measure, and our
repeated-inference protocol provides exactly the samples such a scheme
would aggregate. \textbf{Semantic entropy}~\cite{farquhar2024semantic}
clusters generations by meaning and estimates uncertainty in semantic
space to flag confabulations; our semantic self-similarity metric is a
lightweight, deployment-oriented cousin of this idea, trading a full
entropy estimate for an interpretable pairwise-agreement score.
\textbf{Conformal prediction}~\cite{angelopoulos2021conformal} offers
distribution-free guarantees and could wrap a model's repeated outputs
in calibrated prediction sets, converting raw variability into a
formal coverage statement suitable for safety-critical review.

A further promising direction, and one we highlight at the editor's
suggestion, is to model the higher-order uncertainty we observe with a
\textbf{type-3 fuzzy logic / ANFIS} layer. Type-3 fuzzy
systems extend type-1 and interval type-2 systems with an additional
level of membership-function uncertainty, which makes them well suited
to representing the ``uncertainty about uncertainty'' that arises when a
model returns many semantically related but non-identical answers across
runs~\cite{mohammadzadeh2020type3, cao2021type3}. Concretely, an
adaptive neuro-fuzzy inference system (ANFIS) with type-3 antecedents
could act as an uncertainty-aware aggregator over the $N$ repeated
generations, mapping features of the run set (e.g., pairwise semantic
similarity, judge-score spread, response-length dispersion) to a
calibrated stability/confidence score that gates whether an answer is
surfaced, escalated, or withheld. We leave a full implementation and
evaluation of such a type-3 ANFIS post-processing layer to future work,
but note it as a principled route from \emph{measuring} reproducibility
to \emph{actively managing} it in deployment.

\paragraph{Inference environment.}
All results were obtained via the Ollama local inference server on a
single workstation.
Quantization level, hardware configuration, and batch size can affect
generation quality and throughput; API-hosted counterparts of the same
model families may exhibit different variance profiles.

% ---- 6. Conclusion ------------------------------------------
\section*{Conclusion}

We have presented and applied a reproducibility-first evaluation
framework for small open-weight LLMs in clinical question answering,
covering 1,500 responses from three models evaluated on 50 MedQuAD
questions with 10 runs each.
The results lead to three concrete takeaways for practitioners.

First, \textbf{low temperature does not buy content reproducibility}.
Even at $T{=}0.2$, self-agreement rates hover between 0.12 and 0.20,
and 87--97\% of all outputs per model are unique.
Any medical related usage must account for this variance through
ensembling, gating, or human review.

Second, \textbf{at these parameter counts, clinical fine-tuning does
not compensate for a large capacity gap}.
Gemma 3 12B (general-purpose, 12B) outperforms MedGemma 1.5 4B
(clinically fine-tuned, 4B) on every dimension: judge score, lexical
quality, and reproducibility.
However, this comparison confounds fine-tuning with a 3$\times$ scale
difference; a same-size comparison (e.g., Gemma 3 4B vs.\
MedGemma 1.5 4B) is needed to isolate the fine-tuning effect.
Model selection should be empirically driven, not label-driven.

Third, \textbf{quality and reproducibility partially trade off}.
Llama 3.1 8B maximizes lexical quality and speed but minimizes
consistency; Gemma 3 12B offers the best balance of quality and
reproducibility at the cost of throughput.
The right choice depends on whether the deployment context prioritizes
answer fidelity, consistency, or latency.

Our open-source pipeline makes this three-dimensional evaluation
accessible to any team with a commodity workstation, and we invite
the clinical NLP community to adopt reproducibility metrics as a
standard alongside accuracy in LLM benchmarking.

\paragraph{Future work.}
Planned extensions include additional models (Phi-3, Mistral-Medical,
Qwen-Med), a controlled comparison of same-size general-purpose and
clinically fine-tuned models (e.g., Gemma 3 4B vs.\
MedGemma 1.5 4B) to isolate the effect of domain adaptation from
model scale, evaluation on USMLE and MedQA benchmarks, ensemble
decoding strategies to improve reproducibility, and a calibration
study of the LLM judge against expert clinician ratings on a
real-world EHR-derived question set
(e.g., MIMIC-IV~\cite{johnson2023mimic}).

% ---- Supporting information ---------------------------------
\nolinenumbers
\section*{Supporting information}

\paragraph*{S1 File.}
\label{S1_File}
{\bf Sampled evaluation question set.} The exact 50 MedQuAD questions
used for evaluation, regenerable with
\texttt{clinical-eval export-sample -{}-sample-random 50 -{}-sample-seed 42}.

% ---- References ---------------------------------------------
\begin{thebibliography}{99}

\bibitem{abacha2019mqa}
A.~Ben Abacha and D.~Demner-Fushman.
\newblock A question entailment approach to question answering.
\newblock \emph{BMC Bioinformatics}, 20(1):511, 2019.

\bibitem{dubey2024llama3}
A.~Dubey et~al.
\newblock The Llama 3 herd of models.
\newblock \emph{arXiv preprint arXiv:2407.21783}, 2024.

\bibitem{jin2021disease}
D.~Jin et~al.
\newblock What disease does this patient have? A large-scale open domain
question answering dataset from medical exams.
\newblock \emph{Applied Sciences}, 11(14):6421, 2021.

\bibitem{jin2019pubmedqa}
Q.~Jin et~al.
\newblock PubMedQA: A biomedical research question answering dataset.
\newblock In \emph{Proc.\ EMNLP}, 2019.

\bibitem{johnson2023mimic}
A.~E.~W. Johnson et~al.
\newblock MIMIC-IV, a freely accessible electronic health record dataset.
\newblock \emph{Scientific Data}, 10(1):1, 2023.

\bibitem{lin2004rouge}
C.-Y. Lin.
\newblock ROUGE: A package for automatic evaluation of summaries.
\newblock In \emph{Text Summarization Branches Out (ACL Workshop)}, 2004.

\bibitem{medgemma2024}
{Google DeepMind / medaibase contributors}.
\newblock MedGemma: Clinically fine-tuned Gemma models.
\newblock \url{https://huggingface.co/medaibase}, 2024.

\bibitem{ollama}
{Ollama Contributors}.
\newblock Ollama: Run large language models locally.
\newblock \url{https://ollama.com}, 2024.

\bibitem{papineni2002bleu}
K.~Papineni et~al.
\newblock BLEU: A method for automatic evaluation of machine translation.
\newblock In \emph{Proc.\ ACL}, 2002.

\bibitem{singhal2023large}
K.~Singhal et~al.
\newblock Large language models encode clinical knowledge.
\newblock \emph{Nature}, 620(7972):172--180, 2023.

\bibitem{team2024gemma}
{Gemma Team, Google DeepMind}.
\newblock Gemma 3 technical report.
\newblock \emph{arXiv preprint arXiv:2503.19786}, 2025.

\bibitem{zhang2020bertscore}
T.~Zhang et~al.
\newblock BERTScore: Evaluating text generation with BERT.
\newblock In \emph{Proc.\ ICLR}, 2020.

\bibitem{zheng2023judging}
L.~Zheng et~al.
\newblock Judging LLM-as-a-judge with MT-Bench and chatbot arena.
\newblock In \emph{NeurIPS}, 2023.

\bibitem{wang2022selfconsistency}
X.~Wang et~al.
\newblock Self-consistency improves chain of thought reasoning in language
models.
\newblock In \emph{Proc.\ ICLR}, 2023.

\bibitem{chen2023universal}
X.~Chen et~al.
\newblock Universal self-consistency for large language model generation.
\newblock \emph{arXiv preprint arXiv:2311.17311}, 2023.

\bibitem{kadavath2022calibration}
S.~Kadavath et~al.
\newblock Language models (mostly) know what they know.
\newblock \emph{arXiv preprint arXiv:2207.05221}, 2022.

\bibitem{nori2023gpt4med}
H.~Nori et~al.
\newblock Can generalist foundation models outcompete special-purpose tuning?
Case study in medicine.
\newblock \emph{arXiv preprint arXiv:2311.16452}, 2023.

\bibitem{farquhar2024semantic}
S.~Farquhar, J.~Kossen, L.~Kuhn, and Y.~Gal.
\newblock Detecting hallucinations in large language models using semantic
entropy.
\newblock \emph{Nature}, 630:625--630, 2024.

\bibitem{angelopoulos2021conformal}
A.~N. Angelopoulos and S.~Bates.
\newblock A gentle introduction to conformal prediction and distribution-free
uncertainty quantification.
\newblock \emph{arXiv preprint arXiv:2107.07511}, 2021.

\bibitem{mohammadzadeh2020type3}
A.~Mohammadzadeh, M.~H. Sabzalian, and W.~Zhang.
\newblock An interval type-3 fuzzy system and a new online fractional-order
learning algorithm: Theory and practice.
\newblock \emph{IEEE Transactions on Fuzzy Systems}, 28(9):1940--1950, 2020.

\bibitem{cao2021type3}
Y.~Cao, A.~Raise, A.~Mohammadzadeh, et~al.
\newblock Deep learned recurrent type-3 fuzzy system: Application for renewable
energy modeling/prediction.
\newblock \emph{Energy Reports}, 7:8115--8127, 2021.

\bibitem{suarezlledo2021misinfo}
V.~Suarez-Lledo and J.~Alvarez-Galvez.
\newblock Prevalence of health misinformation on social media: Systematic
review.
\newblock \emph{Journal of Medical Internet Research}, 23(1):e17187, 2021.

\bibitem{sager2021reddit}
M.~A. Sager et~al.
\newblock Identifying and responding to health misinformation on Reddit
dermatology forums with artificially intelligent bots using natural language
processing: Design and evaluation study.
\newblock \emph{JMIR Dermatology}, 4(2):e20975, 2021.

\end{thebibliography}

\end{document}
